\chapter{Build: An Order-Book Model, End to End}\label{ch:ml:build-an-order-book-model-end-to-end}

The book's last model is small on purpose: a forecast of the next second of the mid price, trained on the simulator's sessions,
validated without leaks, compiled into a few hundred nanoseconds, and traded in the same simulator, where every microsecond it
takes is a microsecond of stale decisions. Every step reuses a component built earlier in the book, and the whole run is one
reproducible graph. The chapter ends on the question the title of its problem asks: how much is a microsecond worth to this model?
The answer is not what the title suggests. The model earns \$21.73 per twenty-minute session over six test sessions, four times
what the same quoting earns without it; its edge is unchanged up to 50 milliseconds of decision time and falls to zero at about
106 milliseconds. In a market whose top of book changes four times a second, microseconds are worth nothing and a tenth of a
second is worth everything; what the latency costs depends on how fast the market moves, not on how fast the code runs.

\section{The specification}

The model, written as the \omterm{def:ml:the-machine-learning-team:handoff}{handoff specification} of chapter~28 would have it:
\begin{itemize}
\item \emph{Market.} Book~10's exchange simulator, one venue and one instrument at \$100 with a one-cent tick, default fees (a
rebate of \$0.002 a share for providing liquidity, \$0.003 for taking it), Book~7's \texttt{firm\_tape} flow as background, the
agent behind Book~10's latency model (20 microseconds each way).
\item \emph{Decision.} At every change of the top of book the agent sees: compute ten features, forecast the mid change over the
next second (in ticks), and act.
\item \emph{Action.} Quote one lot (100 shares) at the touch on each side its position allows (at most one lot long or short),
except the side the forecast says is about to be run over: the bid is withdrawn when the forecast is below $-\theta$, the ask when
it is above $\theta$. A quote whose price is still the touch is left alone, since its place in the queue is worth keeping. A
position held 30 seconds is closed at the market.
\item \emph{Sessions.} Twenty minutes each, on disjoint seeds: eight for training, three for validation, six for testing.
\item \emph{Budget and monitoring.} The decision time is the compiled model's measured latency; chapter~27's input monitors run on
the test sessions.
\end{itemize}

The strategy is a market maker that uses the model to avoid being picked off, the classic use of a short-horizon forecast by a
liquidity provider (Cartea, Jaimungal and Penalva treat its optimal form). A taker version was tried first and abandoned: in this
market the mid moves by a tick or more within a second about five times in a hundred, and crossing a one-tick spread and paying
the taker fee twice costs 1.6 ticks a round trip, more than any forecast here is worth.

The whole run is one \texttt{firm.workflow} graph (\cref{lst:ml:lob:graph}): two data stages, purged cross-validation, the fit
and its export, the TCN challenger, the choice of threshold on the validation sessions, and the test sessions at every latency.
Each stage is keyed by its code, parameters, seed and inputs, so a change anywhere reruns exactly what depends on it.

\section{Data, labels and validation}

The research sessions use the same agent class as production, with trading switched off: it records its feed as the events of
chapter~24's \omterm{def:ml:data-and-feature-stores:store}{feature store} (each change of the top of book, with its order-flow imbalance increment, and each trade with its
aggressor's sign) and computes its features online at every book change (\cref{lst:ml:lob:onbook}). The ten features are the
imbalance at the touch, the spread, the order-flow imbalance of Cont, Kukanov and Stoikov over one and five seconds, signed traded
volume over one and five seconds, the number of trades over five seconds, the mid change over one and five seconds, and the number
of book updates over one second. The same definitions computed offline over the recorded events equal the online values exactly
at every decision (the acceptance test of \texttt{firm.lobmodel}); training and serving cannot drift apart because they are the
same code.

The label is chapter~2's \omterm{def:ml:targets-labels-and-sample-weights:fixed}{fixed-horizon label} in event time (\cref{lst:ml:lob:data}): the mid change, in ticks, from the last
event known at the decision to the last one a second later, computed by \texttt{firm.labeling} on the per-event mid changes.
Most labels are zero (89.4\%): in a second the mid usually does not move. The eight training sessions give 51,924 decisions and the
three validation sessions 33,457; activity differs a lot from session to session.

Validation respects time. Each label spans one second, so neighbouring rows share their future; chapter~3's purged five-fold
cross-validation (\texttt{firm.cvsplit}, with a one-second embargo) removes from each training fold the rows whose labels overlap
the test fold. Three LightGBM configurations of 150 trees with 7, 15 and 31 leaves score mean fold information coefficients of
0.558, 0.555 and 0.548; the smallest wins. Refitted on all training sessions, it has an information coefficient of 0.625 in sample
and 0.522 on the validation sessions.

The challenger is chapter~8's \omterm{def:ml:sequence-models-on-order-books:tcn}{temporal convolutional network} (\texttt{firm.lobseq}, 12,867 parameters) on windows of the last 16
decisions' features, trained on seven training sessions and stopped early on the eighth. Its validation information coefficient is
0.452. The forest wins on accuracy, and, as the next section measures, it is also the model that can be made fast; Kolm, Turiel and
Westray found similarly that off-the-shelf networks on stationary order-flow inputs beat networks on raw order-book states.

\section{The model and its compiled form}

The chosen forest has 150 trees of six splits each. Exported to flat arrays by chapter~26's \texttt{firm.mlinfer}, it reproduces
LightGBM's predictions bit for bit on the validation rows; so do the generated C++20 branches, the C++ loop over the arrays and the
Rust kernel, on 200 shared test vectors (\cref{lst:ml:lob:cpp}, \cref{lst:ml:lob:rust}). The agent itself calls a NumPy walk of the
same arrays that adds the leaves in the same order (\cref{lst:ml:lob:predict}).

\begin{table}[htbp]
\centering\small
\begin{tabular}{@{}lrrr@{}}
\toprule
serving path & median & 99th percentile & 99.9th percentile\\
\midrule
LightGBM's Python call (one row) & 263 $\mu$s & 560 $\mu$s & 869 $\mu$s\\
TCN forward pass in PyTorch (one window) & 166 $\mu$s & 460 $\mu$s & 718 $\mu$s\\
flat forest in Python and NumPy & 37 $\mu$s & 98 $\mu$s & 172 $\mu$s\\
flat forest, C++20 loop & 1.1 $\mu$s & 2.5 $\mu$s & 9.1 $\mu$s\\
flat forest, generated C++20 branches & 153 ns & 383 ns & 1.1 $\mu$s\\
\bottomrule
\end{tabular}
\caption{Latency of one prediction, measured once on a laptop (Intel Core Ultra 7 155H, WSL2, a shared machine without isolated
cores; the clock's own overhead is about 15 nanoseconds). Data: \texttt{bench\_lobmodel.py},
\texttt{measured\_latency.csv}.}\label{tab:ml:lob:latency}
\end{table}

\Cref{tab:ml:lob:latency} gives the budget's other side. The generated branches answer in 153 nanoseconds at the median, about
1,700 times faster than LightGBM's Python call; the TCN, less accurate, would also be a thousand times slower than the compiled
forest. The trading runs use these measurements as the agent's declared decision time: 150 nanoseconds for the compiled model,
263 microseconds for the library call.

\section{Trading it in the simulator}

\begin{definition}[Decision staleness]\label{def:ml:build-an-order-book-model-end-to-end:staleness}
The \emph{decision staleness}\index{decision staleness} of an order is the time from the market event its decision is based on to
the moment the order leaves the trader: the market-data latency, any wait while earlier decisions are still being computed, and the
decision time itself. With decisions triggered at a rate $\lambda$ and each taking $L$, the wait stays bounded only while
$\lambda L<1$.
\end{definition}

The agent's decision (\cref{lst:ml:lob:act}) runs at every book change it sees. Its threshold is chosen on the validation sessions
at the compiled model's latency: mean net P\&L per session of \$43.4, \$53.9, \$48.0 and \$50.1 for $\theta$ of 0.05, 0.1, 0.15 and
0.2 ticks, so $\theta=0.1$. On the six test sessions, at 150 nanoseconds, the model earns \$21.73 per session (standard error
\$3.90), positive in all six, from 34.3 entries per session; its net exchange fees are a credit of \$6.73. Its entries' mark-out one second
later is 0.405 ticks in its favour. The same quoting without a model (never withdrawing a side) earns \$5.47 (standard error
\$2.22) from 58.7 entries, with a one-second mark-out of 0.078 ticks: it is filled more often and mostly by orders that know
better. The model's value is in the fills it declines.

The test sessions also run chapter~27's input monitor, each session against the training sessions. The largest \omterm{def:ml:monitoring-and-retraining:psi}{population
stability index} of a session ranges from 0.14 to 1.60, far above any rule of thumb, and in every session it comes from a feature
that counts events (updates in the last second, trades in the last five), which moves with the session's activity: the test
sessions see between 2.8 and 6.4 book changes a second. The four price features stay between 0.02 and 0.13. The model is profitable
in all six sessions. A monitor on count features needs its reference to span the range of activity the model will meet, or
features normalised by activity, or it pages every quiet afternoon.

\section{What the latency costs}

The test sessions are traded again with the decision time set to each point of a grid from zero to 150 milliseconds
(\cref{tab:ml:lob:results}, \cref{fig:ml:lob:pnl}). Everything else is identical: the sessions, the model, the threshold.

\begin{table}[htbp]
\centering\small
\begin{tabular}{@{}rrrrr@{}}
\toprule
decision time & median staleness & net P\&L per session (\$) & 1-s mark-out (ticks) & sessions positive\\
\midrule
0 & 0.02 ms & 21.73 & 0.405 & 6 of 6\\
150 ns & 0.02 ms & 21.73 & 0.405 & 6 of 6\\
263 $\mu$s & 0.28 ms & 21.73 & 0.405 & 6 of 6\\
1 ms & 1.02 ms & 21.23 & 0.399 & 6 of 6\\
10 ms & 10.0 ms & 21.72 & 0.405 & 6 of 6\\
30 ms & 30.0 ms & 20.63 & 0.394 & 6 of 6\\
50 ms & 50.0 ms & 20.52 & 0.417 & 6 of 6\\
70 ms & 70.0 ms & 12.40 & 0.353 & 6 of 6\\
100 ms & 154 ms & 1.60 & 0.245 & 3 of 6\\
150 ms & 802 ms & $-12.27$ & 0.148 & 2 of 6\\
no model & 0.02 ms & 5.47 & 0.078 & \\
\bottomrule
\end{tabular}
\caption{The model on six test sessions at each decision time: median \omterm{def:ml:build-an-order-book-model-end-to-end:staleness}{decision staleness}, mean net P\&L per twenty-minute session,
mean one-second mark-out of entries, and the sessions with a positive P\&L; the last row quotes without a model. Data:
\texttt{ml\_lobmodel.latency\_table}.}\label{tab:ml:lob:results}
\end{table}

\begin{omfigure}
\begin{tikzpicture}
\begin{axis}[width=11.5cm,height=6cm,xmode=log,xlabel={decision time (ms)},ylabel={net P\&L per session (\$)},xmin=1e-4,xmax=300,
  ymin=-22,ymax=30,tick label style={font=\small},label style={font=\small},grid=major,grid style={gray!25},legend pos=south west,
  legend cell align=left,legend style={font=\small},table/col sep=comma]
\addplot[name path=hi,draw=none,forget plot] table[x=ms,y=hi]{figdata/ml/29-build-an-order-book-model-end-to-end/latency.csv};
\addplot[name path=lo,draw=none,forget plot] table[x=ms,y=lo]{figdata/ml/29-build-an-order-book-model-end-to-end/latency.csv};
\addplot[omProp!25,forget plot] fill between[of=hi and lo];
\addplot[omProp,very thick,mark=*,mark size=1.5pt] table[x=ms,y=pnl]{figdata/ml/29-build-an-order-book-model-end-to-end/latency.csv};
\addplot[omDef,thick,dashed] table[x=ms,y=base]{figdata/ml/29-build-an-order-book-model-end-to-end/latency.csv};
\addplot[black,thin] coordinates {(1e-4,0) (300,0)};
\legend{the model ($\pm$ one standard error),quoting without a model}
\end{axis}
\end{tikzpicture}

\omcaption{Mean net P\&L per twenty-minute test session against the decision time, from the compiled model's 150 nanoseconds to
150 milliseconds (log scale; zero gives the same as 150 nanoseconds), with a band of one standard error over the six sessions;
dashed: the same quoting without a model. Data: \texttt{ml\_lobmodel.latency\_table}.}\label{fig:ml:lob:pnl}
\end{omfigure}

Up to 50 milliseconds nothing changes: the P\&L stays within its standard error of \$21.73, and the mark-out within 0.02 ticks of
0.405. Beyond, the edge goes quickly. At 70 milliseconds the P\&L is \$12.40; at 100 milliseconds it is \$1.60, positive in
only three sessions; at 150 milliseconds it is $-\$12.27$. By linear interpolation on the grid, the P\&L falls to zero at about 106
milliseconds, and to what quoting without a model earns at about 89 milliseconds: past that, the model is worse than no model.

Two mechanisms produce the cliff, and the staleness column separates them. Up to 70 milliseconds the median staleness equals the
decision time: each decision is late by exactly that much, and a quote withdrawn 50 milliseconds late is rarely hit in those 50
milliseconds, because the book changes only about four times a second. At 100 milliseconds the median staleness is 154
milliseconds: in the busiest bursts, decisions arrive faster than one every 100 milliseconds, $\lambda L$ passes one, and they queue
behind each other. At 150 milliseconds the median staleness is 802 milliseconds, and the agent quotes on a book that is almost a
second old. A model that is slow on average is slower still exactly when the market is busy, which is when its forecasts matter.

The one-second mark-out tells the same story from the fills' side, falling from 0.405 to 0.148 ticks; it stays positive, because
even a late withdrawal avoids some of the worst fills, but not by enough to pay for the rest.

Why do microseconds not matter here? Because the time that counts is measured in the market's events, not in seconds. Kolm,
Turiel and Westray found that the effective horizon of their stock-specific forecasts was about two average price changes. In this
market the book changes about four times a second and no one else is racing for the same quotes: the background flow does not
react to the agent. On a real exchange the top of a liquid stock changes far more often than four times a second, and the agent's quotes are
picked off by other fast traders: Aquilina, Budish and O'Neill, from exchange message data, find latency-arbitrage races about
once a minute per FTSE~100 stock, with the modal race lasting 5 to 10 millionths of a second. The same cliff would sit
correspondingly closer to zero, and the microseconds that bought nothing here would be the whole business. The simulator
answers the question it was built to answer, the cost of staleness in its own event time; carrying the answer to a real market
means rescaling it by that market's event rate and adding its competitors.

\begin{method}[Building a short-horizon model that trades]\label{met:ml:lob:method}
\begin{enumerate}
\item Specify the market, the decision, the action and the budget before the model.
\item Record research data with the production agent, so features are the same code in both; test that offline equals online.
\item Validate with purged folds; choose the threshold on sessions the model was not fitted on; test on others.
\item Compile the chosen model and check parity on shared vectors; measure its latency and use the measurement.
\item Trade it at a grid of decision times, and report P\&L, mark-outs and staleness against the time, next to a no-model baseline.
\item Read latency in the market's event time: the edge's horizon is a number of price changes, not of microseconds.
\end{enumerate}
\end{method}

\section{Tutorial: microseconds are money}\label{tut:ml:build-an-order-book-model-end-to-end:tut}

\begin{tutorial}
\textbf{Goal.} Run the whole pipeline, compile the model, trade it at every decision time, and find where its edge ends.
\textbf{End state:} \cref{tab:ml:lob:latency}, \cref{tab:ml:lob:results}, \cref{fig:ml:lob:pnl}.
\begin{tutsteps}
\item \textbf{The graph.}
\omcode{code/firm/lobmodel/firm_lobmodel.py}{344}{358}{The chapter as one firm.workflow graph.}\label{lst:ml:lob:graph}
\item \textbf{Features at every book change.}
\omcode{code/firm/lobmodel/firm_lobmodel.py}{92}{111}{The agent records an event, computes its features online, forecasts and declares its decision time.}\label{lst:ml:lob:onbook}
\item \textbf{Labels.}
\omcode{code/firm/lobmodel/firm_lobmodel.py}{167}{177}{Features offline and the fixed-horizon label in event time.}\label{lst:ml:lob:data}
\item \textbf{Serving in Python.}
\omcode{code/firm/lobmodel/firm_lobmodel.py}{205}{225}{The flat forest walked for all trees at once, leaves added in order.}\label{lst:ml:lob:predict}
\item \textbf{Serving in C++20.}
\omcode{code/firm/lobmodel/cpp/lobmodel.hpp}{1}{23}{The C++20 serving path: firm.mlinfer's loop and the generated branches.}\label{lst:ml:lob:cpp}
\item \textbf{Serving in Rust.}
\omcode{code/firm/lobmodel/rust/src/lib.rs}{5}{23}{The Rust serving path on firm.mlinfer's kernel.}\label{lst:ml:lob:rust}
\item \textbf{The decision.}
\omcode{code/firm/lobmodel/firm_lobmodel.py}{113}{130}{Quote each allowed side at the touch unless the forecast says it is about to be run over.}\label{lst:ml:lob:act}
\item \textbf{Run} \texttt{ml\_lobmodel.summary()}, \texttt{latency\_table()}, \texttt{zero\_crossing()}, \texttt{monitors()},
\texttt{fig\_lobmodel.py} (about three minutes on one core), and \texttt{bench\_lobmodel.py} once for the measured latencies.
\end{tutsteps}
\textbf{What to change next.} Add a second, faster quoting agent to the simulator that withdraws on the same signal: the cliff
should move towards the difference between the two agents' decision times.
\end{tutorial}

\section{Build: the order-book model}\label{bld:ml:build-an-order-book-model-end-to-end:lobmodel}

\begin{build}
\bfield{Purpose} A short-horizon model from simulated sessions to trading, with every step reproducible and its serving path in
three languages.
\bfield{Interface} \texttt{FEATURES}, \texttt{HORIZON}, \texttt{DEFAULT}; \texttt{ModelAgent(predict, threshold, hold\_s,
decision\_ns, qty, trade)} (a \texttt{firm.exchsim} agent); \texttt{run\_session}, \texttt{record}, \texttt{events\_array},
\texttt{dataset}, \texttt{report}; \texttt{ForestPredictor}; the stages \texttt{stage\_data}, \texttt{stage\_cv},
\texttt{stage\_fit}, \texttt{stage\_tcn}, \texttt{stage\_threshold}, \texttt{stage\_test}; \texttt{make\_pipeline(cache\_dir,
cfg)}; \texttt{cpp/lobmodel.hpp} and the Rust crate \texttt{firm\_lobmodel}; \texttt{make\_lobmodel\_fixture.py}.
\bfield{Rules} Research and production share the agent and the feature code; sessions of different roles never share a seed; the
decision time is a measured latency; the serving paths agree bit for bit.
\bfield{Acceptance tests} \texttt{code/firm/lobmodel/tests/}: online features equal offline ones at every decision and labels
equal the mid change; the flat forest, the NumPy walk and LightGBM agree on the fixture, and the pipeline rebuilds the fixture's
forest; a longer decision time makes decisions staler. \texttt{cpp/lobmodel\_test.cpp} and \texttt{cargo test}: exact parity on
200 vectors.
\bfield{Stretch} A competing fast agent; features normalised by activity; the TCN compiled with chapter~26's int8 path.
\end{build}

\begin{omsources}
\item P.~N.~Kolm, J.~Turiel and N.~Westray, ``Deep order flow imbalance: extracting alpha at multiple horizons from the limit
order book'', \emph{Mathematical Finance}, 2023.
\item J.~Sirignano and R.~Cont, ``Universal features of price formation in financial markets: perspectives from deep learning'',
\emph{Quantitative Finance}, 2019.
\item R.~Cont, A.~Kukanov and S.~Stoikov, ``The price impact of order book events'', \emph{Journal of Financial Econometrics},
2014.
\item Á.~Cartea, S.~Jaimungal and J.~Penalva, \emph{Algorithmic and High-Frequency Trading}, Cambridge University Press, 2015.
\item M.~Aquilina, E.~Budish and P.~O'Neill, ``Quantifying the high-frequency trading `arms race''', \emph{Quarterly Journal of
Economics}, 2022.
\end{omsources}

\section{Exercises}

\begin{exercise}[$\star$]\label{exo:ml:build-an-order-book-model-end-to-end:1}
Why does the agent record its research data itself, rather than the research using the simulator's full tape?
\end{exercise}

\begin{exercise}[$\star$]\label{exo:ml:build-an-order-book-model-end-to-end:2}
Why was a taker strategy abandoned? Compute its round-trip cost in ticks.
\end{exercise}

\begin{exercise}[$\star$]\label{exo:ml:build-an-order-book-model-end-to-end:3}
The model earns more than quoting without it, yet enters less often. Explain with the mark-outs.
\end{exercise}

\begin{exercise}[$\star\star$]\label{exo:ml:build-an-order-book-model-end-to-end:4}
At 100 milliseconds the median staleness is 154 milliseconds. Explain, with $\lambda L$, and say what a production system would do
about it.
\end{exercise}

\begin{exercise}[$\star\star$]\label{exo:ml:build-an-order-book-model-end-to-end:5}
The test sessions' input monitor reads up to 1.60. Should the model have been stopped? What would you change in the monitor?
\end{exercise}

\begin{exercise}[$\star\star$]\label{exo:ml:build-an-order-book-model-end-to-end:6}
\emph{Find the flaw.} ``The compiled model is 1,700 times faster than LightGBM's Python call, so it will make 1,700 times more
money.''
\end{exercise}

\begin{exercise}[$\star\star\star$]\label{exo:ml:build-an-order-book-model-end-to-end:7}
Where would the zero-crossing sit for a stock whose top of book changes 2,000 times a second, if the edge's horizon is a fixed
number of book changes? What does the estimate leave out?
\end{exercise}

\begin{exercise}[$\star\star\star$]\label{exo:ml:build-an-order-book-model-end-to-end:8}
Write the monitoring specification for this model in production: the monitors, their references, thresholds and owners.
\end{exercise}

\section{Problem: Microseconds Are Money}

\begin{problem}[{Weekend problem --- microseconds are money}]\label{pb:ml:build-an-order-book-model-end-to-end:1}
The chapter's pipeline, model and latency grid.

\medskip\noindent\textbf{Part I --- Specification and data.}
\begin{enumerate}
\item State the market, the decision, the action and the sessions.
\item What are the features, and why do offline and online values agree?
\item How is the label defined, and why are most labels zero?
\item How is the model validated, and what does the cross-validation choose?
\end{enumerate}

\medskip\noindent\textbf{Part II --- The model.}
\begin{enumerate}[resume]
\item How do the forest and the TCN compare?
\item What does the compiled forest cost per prediction, and against what?
\item How is parity checked?
\item How is the threshold chosen?
\end{enumerate}

\medskip\noindent\textbf{Part III --- Trading.}
\begin{enumerate}[resume]
\item What does the model earn at 150 nanoseconds, and against what baseline?
\item What do the mark-outs show?
\item What does the input monitor report, and why?
\item Define \omterm{def:ml:build-an-order-book-model-end-to-end:staleness}{decision staleness} and explain its two parts.
\end{enumerate}

\medskip\noindent\textbf{Part IV --- The verdict.}
\begin{enumerate}[resume]
\item State the \emph{named result}: net P\&L per session and the one-second mark-out as functions of decision latency, and the
latency at which the model's edge falls to zero.
\item Why is the curve flat up to 50 milliseconds?
\item Why does it fall so fast after?
\item What would change on a real exchange?
\item Should the firm spend on the C++ serving path for this market?
\item What would you add to the simulator before trusting the curve?
\item How would you register and monitor this model?
\item In one sentence: what is a microsecond worth?
\end{enumerate}
\end{problem}

\section{Interview questions}

\begin{interviewq}[$\star$ \iqroles{researcher, mle}]\label{iq:ml:build-an-order-book-model-end-to-end:1}
What features would you use to forecast the next second of a stock's mid price?
\end{interviewq}

\begin{interviewq}[$\star\star$ \iqroles{researcher}]\label{iq:ml:build-an-order-book-model-end-to-end:2}
How do you validate a model whose labels overlap in time?
\end{interviewq}

\begin{interviewq}[$\star\star$ \iqroles{mle, dev}]\label{iq:ml:build-an-order-book-model-end-to-end:3}
How do you make sure a model computes the same features in production as in research?
\end{interviewq}

\begin{interviewq}[$\star\star$ \iqroles{researcher, trader}]\label{iq:ml:build-an-order-book-model-end-to-end:4}
A market maker adds a short-horizon forecast. Where does its value show up, and how do you measure it?
\end{interviewq}

\begin{interviewq}[$\star\star$ \iqroles{mle, dev}]\label{iq:ml:build-an-order-book-model-end-to-end:5}
Your model takes 200 microseconds in Python. How do you get it under a microsecond, and how do you know it still gives the same
answers?
\end{interviewq}

\begin{interviewq}[$\star\star\star$ \iqroles{researcher, dev}]\label{iq:ml:build-an-order-book-model-end-to-end:6}
How would you decide how much to spend on latency for a given strategy?
\end{interviewq}
